\subsection{DETERMINANTS OF AN ENDOMORPHISM}

Let M be an A-module with a finite basis of n elements and u an endomorphism of M. The A-module \(\bigwedge^{n}(\mathbf{M})\) is a monogenous free module, that is isomorphic to A (no. 8, Corollary 1 to Theorem 1); \(\bigwedge^{n}(u)\) is an endomorphism of this module and is therefore a homothety \(z \mapsto \lambda z\) of ratio \(\lambda \in A\) determined uniquely (II, §2, no. 3, Proposition 5).

DEFINITION 1. The determinant of an endomorphism \(u\) of a free A-module \(M\) of finite dimension \(n\) (II, §7, no. 2, Corollary to Proposition 3 and Remark 1), denoted by \(\det u\) , is the scalar \(\lambda\) such that \(\bigwedge^{n}(u)\) is the homothety of ratio \(\lambda\) .

By formula (4) of §7, no. 2, det u is the unique scalar such that

\[
u \left(x _ {1}\right) \wedge u \left(x _ {2}\right) \wedge \dots \wedge u \left(x _ {n}\right) = (\det u). x _ {1} \wedge x _ {2} \wedge \dots \wedge x _ {n}\tag{1}
\]

for every sequence \((x_{i})_{1\leqslant i\leqslant n}\) of \(n\) elements of M. If \(\det (u) = 1\) , \(u\) is said to be unimodular.

PROPOSITION 1. (i) If \(u\) and \(v\) are two endomorphisms of a finite-dimensional free A-module M, then

\[
\det (u \circ v) = (\det u) (\det v).\tag{2}
\]

\begin{enumerate}
\def\labelenumi{(\roman{enumi})}
\setcounter{enumi}{1}
\tightlist
\item
  \(\operatorname{det}(1_{\mathbf{M}}) = 1\) ; for every automorphism \(u\) of \(\mathbf{M}\) , \(\operatorname{det} u\) is invertible in \(\mathbf{A}\) and
\end{enumerate}

\[
\det (u ^ {- 1}) = (\det u) ^ {- 1}.\tag{3}
\]

If \(n\) is the dimension of \(\mathbf{M}\) , this follows immediately from the relation \(\bigwedge^{n}(u \circ v) = (\bigwedge^{n}(u)) \circ (\bigwedge^{n}(v))\) (§ 7, no. 2, formula (3)).

Let M be a free A-module with a finite basis \((e_{i})_{1\leqslant i\leqslant n}\) ; given a sequence \((x_{i})_{1\leqslant i\leqslant n}\) of n elements of M, the determinant of this sequence with respect to the given basis \((e_i)\) , denoted by \(\det(x_1, x_2, \ldots, x_n)\) when no confusion can arise over the basis, is the determinant of the endomorphism \(u\) of M such that \(u(e_i) = x_i\) for \(1 \leqslant i \leqslant n\) . Then by formula (1)

\[
x _ {1} \wedge x _ {2} \wedge \dots \wedge x _ {n} = \det (x _ {1}, x _ {2}, \dots , x _ {n}) e _ {1} \wedge e _ {2} \wedge \dots \wedge e _ {n}\tag{4}
\]

and this relation characterizes the mapping \((x_{i})\mapsto\det(x_{1},x_{2},\ldots,x_{n})\) of \(\mathbf{M}^{n}\) into A. It shows that this mapping is an alternating \(n\) -linear form. As, by virtue of §7, no. 4, Proposition 7, the A-module of alternating \(n\) -linear forms is canonically isomorphic to the dual of \(\bigwedge^{n}(\mathbf{M})\) and \(\bigwedge^{n}(\mathbf{M})\) is isomorphic to A, it is seen that every alternating \(n\) -linear form on \(\mathbf{M}^{n}\) can be written

\[
(x _ {1}, \dots , x _ {n}) \mapsto \alpha \det (x _ {1}, x _ {2}, \dots , x _ {n})
\]

for some \(\alpha \in \mathbf{A}\) .

PROPOSITION 2. Let M be a free A-module with a finite basis \((e_i)_{1 \leqslant i \leqslant n}\) and \(v\) an endomorphism of M. For every sequence \((x_i)_{1 \leqslant i \leqslant n}\) of \(n\) elements of M,

\[
\det (v (x _ {1}), \dots , v (x _ {n})) = (\det v) \det (x _ {1}, \dots , x _ {n}).\tag{5}
\]

If \(u\) is the endomorphism of \(M\) such that \(u(e_i) = x_i\) for all \(i\) , then

\[
v (x _ {i}) = (v \circ u) (e _ {i})
\]

and (5) therefore follows from (2).

